\section{Introduction}
\label{sec:intro}

Dense vector embeddings are a standard representation for semantic retrieval; retrieval-augmented generation, for
example, retrieves the passages closest to a query embedding and passes them to a language model~\cite{lewis2020rag}.
Exact nearest-neighbor search is expensive at scale, so practical systems use approximate nearest-neighbor (ANN)
indexes such as Hierarchical Navigable Small World (HNSW) graphs~\cite{malkov2020hnsw}. These indexes increasingly
live inside general database systems: pgvector adds vector types, distance operators and HNSW indexes to
PostgreSQL~\cite{pgvector}, and Milvus, AnalyticDB-V and VBase integrate vector search with structured data
processing~\cite{wang2021milvus,wei2020analyticdbv,zhang2023vbase}.

Real retrieval requests usually restrict candidates by metadata such as a date range, category or access attribute.
Such a \emph{filtered top-$k$ query} (\code{WHERE attr BETWEEN lo AND hi ORDER BY embedding <=> q LIMIT k}) can be
executed in two fundamentally different ways. \emph{Filter-then-rank} uses a B-tree or a sequential scan to
find every qualifying row and computes its distance to the query. It is exact, but its cost grows linearly with the
number of qualifying rows. \emph{Rank-then-filter} lets an HNSW index scan return candidates in approximate distance
order and discards those that violate the predicate. It is fast when most rows qualify, but under a selective filter
most candidates are discarded. With pgvector's default search width, a filter matching 10\% of the rows leaves only
about four of forty candidates~\cite{pgvector}, so the query can return fewer than $k$ rows. Since version 0.8,
pgvector offers \emph{iterative index scans} that keep traversing the graph until enough rows qualify, up to a limit.
When very few rows qualify, such a scan must visit a large part of the graph.

The better strategy therefore depends on \emph{filter selectivity}, the fraction of rows that satisfy the predicate.
Recent studies confirm that no single filtered-ANN strategy dominates across selectivities, and report that pgvector's
optimizer sometimes chooses an approximate index scan where an exact plan would give perfect recall at similar
latency~\cite{li2025attribute,amanbayev2026filtered}. A classical cost model~\cite{selinger1979access} compares plans
by estimated cost only and has no notion of an index scan that may return an \emph{incomplete} answer.

We therefore ask whether a lightweight \emph{external} routing rule can outperform both fixed strategies and
PostgreSQL's default plan choice while guaranteeing retrieval quality on average. The rule uses only the database's
own cardinality estimates and one threshold calibrated offline. The engine and indexes stay unchanged, which makes the
approach complementary to filter-aware index designs such as Filtered-DiskANN~\cite{gollapudi2023filtered} and
ACORN~\cite{patel2024acorn}. We study three questions.
\textbf{RQ1}: How does selectivity affect the latency, Recall@10 and completeness of exact and HNSW search in
PostgreSQL?
\textbf{RQ2}: Does a calibrated threshold router reduce latency relative to fixed strategies and the default planner on
held-out queries, while meeting a mean Recall@10 $\ge0.95$ and $\ge99\%$ complete queries?
\textbf{RQ3}: Where does the router fail, or where does its overhead outweigh its savings?

Our contributions are: (i) a selectivity-aware router (P1) for unmodified PostgreSQL/pgvector that estimates matching
rows from cached \code{pg\_stats} histograms, without any per-query \code{COUNT(*)}, and dispatches each query to exact
or iterative HNSW search using a threshold calibrated under explicit recall and completeness constraints; (ii) a
controlled benchmark of 2{,}700 filtered queries on 100{,}000 public LAION/CLIP vectors, with plan-level verification of
every strategy, ground truth validated against the benchmark's own identifiers, and a held-out evaluation; and
(iii) an analysis, including negative results, of where routing helps, where the default planner returns incomplete
results, and where a single threshold reaches its limits. Code, configuration and all measured results are
public~\cite{repo}. Sections~\ref{sec:related}--\ref{sec:conclusion} cover related work, the method, the setup, the
results, a discussion with limitations, and the conclusion.
